% Part: first-order-logic
% Chapter: beyond
% Section: many-sorted-logic

\documentclass[../../../include/open-logic-section]{subfiles}

\begin{document}

\olfileid{fol}{byd}{msl}

\olsection{Meersoortige logica}

In de eerste-ordelogica lopen variabelen en kwantoren over één
!!{domain}. Vaak is het echter nuttig om meerdere (disjuncte)
!!{domain}s te hebben: bijvoorbeeld !!a{domain} van getallen,
!!a{domain} van meetkundige objecten, !!a{domain} van functies van
getallen naar getallen, !!a{domain} van abelse groepen, enzovoort.

De meersoortige logica biedt zo'n kader. Men begint met een lijst van
``soorten''; de ``soort'' van een object geeft aan tot welk
``!!{domain}'' het geacht wordt te behoren. Vervolgens zijn er voor
elke soort !!{variable}s en kwantoren, en (doorgaans) een
!!{identity} voor elke soort. Functies en relaties zijn eveneens naar
soort ``getypeerd'': de typen bepalen welke soorten objecten zij als
argumenten kunnen aannemen. Verder blijven de gebruikelijke regels
van de eerste-ordelogica behouden, waarbij voor elke soort een versie
van de kwantorregels wordt gegeven.

Om bijvoorbeeld internationale betrekkingen te bestuderen, zouden we
een taal kunnen kiezen met twee soorten objecten: Franse burgers en
Duitse burgers. Voor objecten van de eerste soort zouden we een unaire
relatie ``drinkt wijn'' kunnen hebben; voor objecten van de tweede
soort een andere unaire relatie, ``eet worst''; en daarnaast een
binaire relatie ``vormen samen een echtpaar met verschillende
nationaliteiten''. Die laatste relatie heeft twee argumenten: het
eerste is van de eerste soort en het tweede van de tweede soort. Als
we de variabelen $a$, $b$, $c$ over Franse burgers laten lopen en $x$, $y$, $z$
over Duitse burgers, dan
\[
\lforall[a][\lforall x][(\Atom{\Obj{MarriedTo}}{a,x} \lif
(\Atom{\Obj{DrinksWine}}{a} \lor \lnot \Atom{\Obj{EatsWurst}}{x})]]
\]
beweert dat, als een Franse burger met een Duitser getrouwd is, de
Franse burger wijn drinkt of de Duitser geen worst eet.

De meersoortige logica kan op natuurlijke wijze in de
eerste-ordelogica worden ingebed. Daartoe neemt men alle objecten uit
de meersoortige !!{domain}s samen in één eerste-orde-!!{domain}, gebruikt
men unaire !!{predicate}s om de soorten bij te houden en relativiseert
men de kwantoren. De eerste-ordetaal die bij het voorbeeld hierboven
hoort, zou naast de andere beschreven relaties bijvoorbeeld de unaire
!!{predicate}s ``$\Obj{German}$'' en ``$\Obj{French}$'' bevatten, waarbij de
soortvereisten uit die andere relaties zijn verwijderd. Een
soortgebonden kwantor $\lforall[x][!A]$, waarbij $x$ !!a{variable} van de
Duitse soort is, wordt vertaald als
\[
\lforall[x][(\Atom{\Obj{German}}{x} \lif !A)].
\]
We moeten axioma's toevoegen die verzekeren dat de soorten gescheiden
zijn---bijvoorbeeld
$\lforall[x][\lnot (\Atom{\Obj{German}}{x} \land
  \Atom{\Obj{French}}{x})]$---en ook axioma's die garanderen dat
``drinkt wijn'' alleen geldt voor objecten die aan het predicaat
$\Atom{\Obj{French}}{x}$ voldoen, enzovoort. Met deze conventies en
axioma's kan men gemakkelijk aantonen dat meersoortige !!{sentence}s
in eerste-orde-!!{sentence}s kunnen worden vertaald, en meersoortige
!!{derivation}s in eerste-orde-!!{derivation}s. Ook meersoortige
!!{structure}s kunnen in overeenkomstige eerste-orde-!!{structure}s
worden ``vertaald'', en omgekeerd. Er geldt dus ook een
volledigheidsstelling voor de meersoortige logica.

\end{document}

